\chapter{Implementare și Testare}
\label{chap:ch4}

\par Capitolul de față detaliază implementarea arhitecturii Sentinel și evaluează eficacitatea sistemului printr-o suită exhaustivă de 21 de scenarii de atac. Fiecare componentă este analizată individual, cu accent pe deciziile arhitecturale netriviale, optimizările de performanță și aspectele de siguranță a execuției care fundamentează corectitudinea sistemului.


\section{Mediul de dezvoltare și testare}
\label{sec:dev_env}

\par Dezvoltarea a fost realizată pe un host Ubuntu 22.04 (kernel 6.17.0). Testarea se desfășoară într-o mașină virtuală Vagrant cu Debian 12 (kernel 6.1.0-amd64, cgroup v2 activat implicit), izolând atacurile de mediul de dezvoltare. Mediul virtualizat este o necesitate absolută: scenariile REP-17 și REP-18 implică injectarea de module de kernel corupte, iar scenariul REP-01 utilizează un rootkit real de tip Diamorphine --- rularea acestora pe sistemul gazdă ar duce inevitabil la compromiterea totală a acestuia sau la instabilitate severă (Kernel Panic).

\par Stiva de unelte utilizată:

\begin{itemize}
    \item \textbf{clang 14+} cu ținta \texttt{bpf} pentru compilarea programului eBPF; flag-ul \texttt{-D\_\_TARGET\_ARCH\_x86} selectează macro-urile corecte de acces la registrele procesorului.
    \item \textbf{Go 1.21+} cu biblioteca \texttt{github.com/cilium/ebpf} \cite{CiliumEBPF} pentru implementarea daemonului de monitorizare și a managerului de carantinare.
    \item \textbf{Python 3.12} cu scikit-learn \cite{Pedregosa2011}, FastAPI și uvicorn pentru motorul de inferență ML și expunerea serviciului REST.
    \item \textbf{SQLite 3.40+} în modul WAL (Write-Ahead Logging) pentru persistența concurentă a datelor.
    \item \textbf{Vagrant 2.4+} cu provider VirtualBox pentru orchestrarea mediului de test.
\end{itemize}


\section{Implementarea programului eBPF}
\label{sec:ebpf_impl}

\par Programul eBPF (\texttt{tracer.bpf.c}, 157 de linii) este compilat cu comanda:

\begin{verbatim}
clang -target bpf -O2 -g -D__TARGET_ARCH_x86 \
  -I/usr/include -I/usr/include/x86_64-linux-gnu \
  -c tracer.bpf.c -o tracer.bpf.o
\end{verbatim}

\par Flag-ul \texttt{-O2} este esențial din două motive: verificatorul eBPF impune o limită de 1 milion de instrucțiuni per execuție și respinge programele cu structuri de control prea complexe; în plus, backend-ul LLVM pentru arhitectura BPF necesită un arbore de sintaxă optimizat pentru a emite \textit{opcodes} BPF valide --- codul neoptimizat (\texttt{-O0}) generează frecvent accesări de memorie pe care Verificatorul le respinge ca nesigure.

\par O decizie de implementare importantă privește extragerea sigură a căilor absolute ale fișierelor. Handler-ul \texttt{submit\_event\_with\_path} utilizează secvența:

\begin{verbatim}
struct pt_regs *inner = (struct pt_regs *)PT_REGS_PARM1(ctx);
unsigned long fname_uptr = 0;
bpf_probe_read_kernel(&fname_uptr, sizeof(fname_uptr), &inner->rsi);
bpf_probe_read_user_str(e->fname, sizeof(e->fname),
                        (const char *)fname_uptr);
\end{verbatim}

\par Kprobe-urile sunt atașate la wrapper-urile \texttt{\_\_x64\_sys\_*} care primesc un singur argument --- un pointer la structura \texttt{pt\_regs} cu argumentele originale ale syscall-ului. Accesul direct la argumentele syscall-ului necesită dereferențierea unui nivel suplimentar de \texttt{pt\_regs} (\textit{inner regs}), o idiomă documentată în \cite{Rice2023} și necesară pentru compatibilitatea cu versiunile kernel $\geq 4.17$ pe x86\_64.

\par Mecanismul de drop count este implementat cu \texttt{\_\_sync\_fetch\_and\_add}, operație atomică per-CPU: când ring buffer-ul este plin și rezervarea eșuează, contorul de drop al CPU-ului curent este incrementat. Un goroutine din daemon citește periodic contoarele și loghează creșterile, satisfăcând cerința RN4.


\section{Implementarea daemonului Go}
\label{sec:daemon_impl}

\par Daemonul Go (\texttt{daemon/main.go}, 497 de linii) orchestrează pipeline-ul de detecție. La pornire, daemonul:

\begin{enumerate}
    \item Ridică limita \texttt{RLIMIT\_MEMLOCK} prin \texttt{rlimit.RemoveMemlock()}, necesară pentru alocarea ring buffer-ului de 64 MiB.
    \item Blochează thread-ul OS cu \texttt{runtime.LockOSThread()} --- această operație este mandatorie, deoarece anumite apeluri de sistem eBPF și asocieri de namespace-uri sunt dependente de contextul specific al thread-ului care le inițiază \cite{CiliumEBPF}.
    \item Încarcă colecția eBPF din fișierul \texttt{.bpf.o}, atașează kprobe-urile și deschide reader-ul de ring buffer.
    \item Inserează propriul PID în \texttt{pid\_blacklist} pentru a preveni buclele de auto-instrumentare.
\end{enumerate}

\par Bucla principală citește din ring buffer în mod blocant (\texttt{rb.Read()}), parsează evenimentele prin \texttt{binary.Read} (little-endian, pentru a respecta ordinea octeților x86\_64) și aplică secvențial cele trei straturi de scoring. Operațiunea de inserare în SQLite utilizează o instrucțiune pregătită (\textit{prepared statement}) pentru a elimina overhead-ul de parsare SQL per eveniment.

\par Detectarea rafale (\texttt{burstTracker}) menține o mapă \texttt{windows[key][]int64} cu timestamp-urile apelurilor dintr-o fereastră glisantă. Cheia de burst este:

\begin{equation}
    \text{key}(pid, nr) = (\text{pid} \ll 32) \mid (nr \ll 16) \mid (\text{start\_time} \mathrel{\&} 0\text{xFFFF})
    \label{eq:burst_key}
\end{equation}

\par Incluzând \texttt{start\_time} (câmpul 22 din \texttt{/proc/\{pid\}/stat}, exprimat în \textit{ticks} de ceas) în cheie, sistemul distinge între procese diferite care rulează cu același PID pe parcursul testelor --- o situație frecventă în suita de atacuri, unde procesele au durată scurtă de viață. Utilizarea exclusivă a ultimilor 16 biți din \texttt{start\_time} introduce o probabilitate de coliziune matematic neglijabilă ($\approx 1/65536$), un compromis excelent între precizia detecției și amprenta de memorie.

\par Un goroutine separat curăță intrările expirate la fiecare 30 de secunde pentru a preveni creșterea nelimitată a memoriei. Altul monitorizează contorul de drop la fiecare 10 secunde.


\section{Implementarea scorerului ML}
\label{sec:scorer_impl}

\par Scorerul ML (\texttt{scorer/isolation\_forest.py}, 195 de linii) expune un API FastAPI cu trei endpoint-uri: \texttt{GET /health}, \texttt{POST /score} și \texttt{POST /model/train}. Modelul este încărcat \textit{lazy} la primul apel \texttt{/score}, eliminând latența de inițializare la pornirea serviciului.

\par Faza de inginerie a caracteristicilor (\textit{feature engineering}) transformă telemetria brută într-un vector de 26 de trăsături. Primele 21 constituie un encoding one-hot al syscall-ului (câte un bit per syscall monitorizat, conform indexului \texttt{SYSCALL\_INDEX}). Caracteristicile suplimentare sunt:

\begin{itemize}
    \item $v_{21}$: flag binar $\mathbb{1}[uid = 0]$ --- capturează comportamentul proceselor root.
    \item $v_{22}$: ponderea de risc normalizată $w_{nr}/100$ --- injectează cunoașterea a priori a riscului syscall-ului în algoritmul nesupervizat. Isolation Forest nu poate distinge semantic între un apel inofensiv și unul critic; adăugarea acestei greutăți forțează algoritmul să izoleze spațial evenimentele inerent periculoase.
    \item $v_{23}$: ora din zi normalizată $h/23$ --- capturează activitate nocturnă atipică.
    \item $v_{24}$: flag weekend $\mathbb{1}[\text{weekday} \geq 5]$.
    \item $v_{25}$: $\log_{10}(\text{pid})/6$ --- diferențiază procesele de sistem (PID mic) de procesele spațiului utilizator.
\end{itemize}

\par Modelul Isolation Forest este antrenat pe evenimentele cu scor static $< 30$ din baza de date (care reprezintă comportamentul de bază), cu parametrii $n\_estimators=200$ și $contamination=0.01$. Vectorii sunt standardizați prin \texttt{StandardScaler} \cite{Pedregosa2011} înaintea antrenării, iar scalarea este aplicată identic la inferență, utilizând parametrii din faza de antrenament. Modelul antrenat este persistat prin \texttt{joblib} și poate fi reantrenat cu date actualizate prin \texttt{POST /model/train} fără a opri serviciul.

\par Conversia scorului brut IF (valoarea \texttt{decision\_function}, negativă pentru anomalii, pozitivă pentru puncte normale) în scor de amenințare pe scala 0--100 utilizează transformarea sigmoidă cu coeficientul de creștere 25, care asigură o tranziție rapidă (abruptă) între stările benigne și maligne --- formula~\ref{eq:sigmoid_transform} din Capitolul~\ref{chap:ch2}. Un punct cu \texttt{raw\_score} $= -0.1$ (ușor anormal) obține \texttt{threat\_score} $\approx 92$, în timp ce un punct cu \texttt{raw\_score} $= 0.1$ (ușor normal) obține \texttt{threat\_score} $\approx 8$.


\section{Implementarea managerului de carantinare}
\label{sec:quarantine_impl}

\par Managerul de carantinare (\texttt{quarantine/main.go}, 306 de linii) validează la pornire disponibilitatea cgroup v2 prin verificarea magic number-ului sistemului de fișiere \texttt{/sys/fs/cgroup} (\texttt{0x63677270} = ``cgrp'' în ASCII) și prezența controllerelor \texttt{cpu} și \texttt{memory} în \texttt{/sys/fs/cgroup/cgroup.controllers}. Această verificare este esențială: pe sisteme cu cgroup v1 sau cgroup v2 fără controllerele activate, acțiunile de carantinare ar eșua silențios.

\par Acțiunile de izolare sunt implementate pe trei niveluri de asprime:

\begin{itemize}
    \item \textbf{THROTTLE}: Creează cgroup \texttt{/sys/fs/cgroup/sentinel/quarantine\_\{pid\}} prin \texttt{os.MkdirAll}, scrie PID-ul procesului în \texttt{cgroup.procs} și scrie \texttt{"10000 100000"} în \texttt{cpu.max}. Această configurare permite procesului să utilizeze 10000 microsecunde din fiecare perioadă de 100000 microsecunde, echivalent cu 10\% CPU --- suficient pentru a sugruma operațiunile intensive (ex. \textit{cryptojacking}), fără a opri complet procesul.
    \item \textbf{ISOLATE}: Creează cgroup, scrie limita de memorie de $256 \times 2^{20}$ octeți în \texttt{memory.max}, apoi aplică regulile nftables. Valoarea de 256 MiB a fost aleasă ca suficientă pentru a permite terminarea normală a procesului, dar insuficientă pentru exfiltrarea unor volume mari de date în memorie. Izolarea de rețea utilizează inodul cgroup-ului ca identificator de filtrare în nftables, evitând necesitatea de network namespaces (care nu pot fi aplicate unui proces deja în execuție fără cooperarea sa).
    \item \textbf{FREEZE}: Trimite \texttt{SIGSTOP} procesului (oprire imediată, indiferent de starea actuală), creează cgroup-ul și scrie \texttt{"1"} în \texttt{cgroup.freeze}. Combinația SIGSTOP + \texttt{cgroup.freeze} este critică: chiar dacă un atacator extern ar încerca să reia procesul printr-un semnal \texttt{SIGCONT}, limitarea impusă la nivel de kernel de cgroup previne planificarea (\textit{scheduling}) procesului de către procesor.
\end{itemize}

\par Izolarea de rețea pentru starea ISOLATE utilizează inodul cgroup-ului ca identificator de filtrare în nftables. Inodul se obține prin \texttt{syscall.Stat} pe calea cgroup-ului; nftables suportă filtrarea \texttt{meta cgroup} prin inod în kernel $\geq 5.10$. Această abordare este net superioară izolării prin \textit{network namespaces}, deoarece poate fi aplicată ``la cald'', fără a necesita repornirea sau cooperarea procesului malițios.

\par Auto-dezghețarea pentru starea FREEZE este gestionată de un goroutine \texttt{unfreezeWatcher} care verifică la fiecare 5 secunde procesele cu deadline expirat. Dezghețarea constă în scrierea \texttt{"0"} în \texttt{cgroup.freeze}, trimiterea \texttt{SIGCONT} și eliminarea procesului din cgroup. Intervalul implicit de 5 minute este configurabil la pornire prin flag-ul \texttt{--unfreeze-after}.


\section{Suita de teste de atac}
\label{sec:attack_suite}

\subsection{Metodologia de testare}
\label{subsec:test_method}

\par Fiecare scenariu de atac este definit printr-un fișier \texttt{attack.toml} cu patru câmpuri: \texttt{id}, \texttt{name}, \texttt{alert\_type} (tipul de alertă așteptat) și \texttt{expected\_state} (starea minimă așteptată). Scriptul \texttt{run.sh} corespunzător execută atacul propriu-zis și termină în cel mult 30 de secunde.

\par Orchestratorul Python \texttt{run\_sentinel.py} automatizează ciclul de testare, garantând izolarea mediului între teste prin golirea jurnalurilor WAL și restaurarea bazei de date la starea de zero:

\begin{enumerate}
    \item Resetare baza de date (\texttt{DELETE FROM events; PRAGMA wal\_checkpoint}).
    \item Execuție \texttt{run.sh} în mașina virtuală prin \texttt{vagrant ssh}.
    \item Așteptare flux de evenimente ($t_{sleep}$ configurabil per scenariu).
    \item Extragere snapshot baza de date din VM prin SCP.
    \item Verificare: există cel puțin un eveniment kprobe (nu sintetic) cu tipul și scorul corecte?
\end{enumerate}

\par Un test trece dacă și numai dacă există cel puțin un eveniment cu sursă \texttt{"kprobe"} (câmpul \texttt{detail.source}) cu \texttt{alert\_type} corespunzător și \texttt{score} $\geq$ pragul minim. Cerința de evenimente reale (kprobe) distinge testele cu detecție reală de eventualele injecții sintetice de testare.

\subsection{Scenariile de bază (REP-01 -- REP-12)}
\label{subsec:basic_attacks}

\par Setul de bază cuprinde 11 scenarii, prezentate în Tabelul~\ref{tab:basic_attacks}, alese să acopere categoriile principale de amenințări identificate de Cozzi et al. \cite{Cozzi2018}.

\begin{table}[htbp]
\begin{center}
\begin{tabular}{|p{45pt}|p{110pt}|p{75pt}|p{55pt}|}
\hline
\textbf{ID} & \textbf{Scenariu} & \textbf{Tip alertă} & \textbf{Stare} \\
\hline\hline
REP-01 & Rootkit Diamorphine                   & HiddenPID           & FREEZE   \\
\hline
REP-02 & Escaladare privilegii (pkexec/setuid) & PrivilegeEscalation & WARN     \\
\hline
REP-04 & Exfiltrare date (netcat)              & NetworkAnomaly      & THROTTLE \\
\hline
REP-05 & Webshell (execve din web server)      & ParentChildAnomaly  & WARN     \\
\hline
REP-06 & Persistență prin cron                & LSMFileWrite        & WARN     \\
\hline
REP-07 & Mapare memorie RWX (mmap)             & MemoryInjection     & WARN     \\
\hline
REP-08 & Manipulare fișiere log               & LogTamper           & WARN     \\
\hline
REP-09 & Injecție bibliotecă LD\_PRELOAD       & LSMFileWrite        & WARN     \\
\hline
REP-10 & Fork bomb                            & SyscallAnomaly      & WATCH    \\
\hline
REP-11 & Injecție prin ptrace                 & ProcessInjection    & FREEZE   \\
\hline
REP-12 & Brute-force SSH                      & NetworkAnomaly      & THROTTLE \\
\hline
\end{tabular}
\end{center}
\caption{Scenariile de atac din setul de bază și rezultatele așteptate}
\label{tab:basic_attacks}
\end{table}

\par REP-01 (rootkit Diamorphine) este implementat ca un modul de kernel inofensiv care nu produce efecte dăunătoare, dar este compilat și instalat prin \texttt{insmod}. Apelul \texttt{init\_module} sau \texttt{finit\_module} produce un scor de 95, declanșând imediat starea FREEZE. Acesta este scenariul cu cel mai scurt timp de detecție din suita de teste: evenimentul apare în ring buffer în milisecunde de la execuția \texttt{insmod}, înainte ca rootkit-ul să își poată ascunde prezența în kernel.

\par REP-11 (injecție ptrace) simulează tehnica clasică de deturnare a memoriei unui proces activ printr-o serie de apeluri \texttt{ptrace(PTRACE\_ATTACH, ...)}, \texttt{ptrace(PTRACE\_POKETEXT, ...)} și \texttt{ptrace(PTRACE\_DETACH, ...)}. Syscall-ul \texttt{ptrace} are ponderea de risc 75; amplificat cu 1.3 pentru root ($97.5$, clamped la 95), produce starea FREEZE.

\par REP-12 (brute-force SSH) generează conexiuni repetate la portul 22 printr-un program C care ciclează \texttt{socket()} + \texttt{connect()} de câteva sute de ori pe secundă. Stratul de burst detection detectează depășirea pragului de 5 apeluri \texttt{connect} în 3 secunde și emite un alert cu scor 75, producând starea THROTTLE.

\subsection{Scenariile avansate (REP-13 -- REP-21)}
\label{subsec:advanced_attacks}

\par Setul avansat cuprinde 9 scenarii cu tehnici mai sofisticate, prezentate în Tabelul~\ref{tab:advanced_attacks}.

\begin{table}[htbp]
\begin{center}
\begin{tabular}{|p{45pt}|p{110pt}|p{75pt}|p{55pt}|}
\hline
\textbf{ID} & \textbf{Scenariu} & \textbf{Tip alertă} & \textbf{Stare} \\
\hline\hline
REP-13 & Execuție fără fișier (memfd)           & MemoryInjection     & WARN     \\
\hline
REP-14 & Process hollowing (ptrace)             & ProcessInjection    & FREEZE   \\
\hline
REP-15 & Reverse shell                          & NetworkAnomaly      & THROTTLE \\
\hline
REP-16 & Recoltare credențiale (/etc/shadow)   & LSMFileWrite        & WARN     \\
\hline
REP-17 & Persistență finit\_module              & HiddenPID           & FREEZE   \\
\hline
REP-18 & Hook syscall (modul kernel)            & HiddenPID           & FREEZE   \\
\hline
REP-19 & Rootkit eBPF (stil TripleCross)        & ProcessInjection    & FREEZE   \\
\hline
REP-20 & Furt memorie proces                   & ProcessInjection    & FREEZE   \\
\hline
REP-21 & Tentativă VM escape                   & PrivilegeEscalation & ISOLATE  \\
\hline
\end{tabular}
\end{center}
\caption{Scenariile de atac din setul avansat și rezultatele așteptate}
\label{tab:advanced_attacks}
\end{table}

\par REP-13 (execuție fără fișier pe disc) simulează tehnica \textit{fileless malware}: codul este copiat într-o zonă de memorie anonimă creată cu \texttt{mmap(MAP\_ANONYMOUS)}, marcată executabilă cu \texttt{mprotect(PROT\_EXEC)}, și apelată direct. Combinația apelurilor \texttt{mmap} (pentru alocare) și \texttt{mprotect} (pentru acordarea drepturilor de execuție zonelor anonime) este detectată cu precizie de motorul logic, producând scoruri de 40, respectiv 45, și rezultând în starea WARN.

\par REP-19 (rootkit eBPF stil TripleCross) reproduce tehnica publicată de Pat Hogan (2022) \cite{Hogan2022}, în care un program eBPF malițios se atașează la apeluri de sistem pentru a modifica comportamentul proceselor. Implementarea din test utilizează un program eBPF inofensiv, dar apelează \texttt{ptrace} pentru a se atașa la un proces țintă, generând scorul de 95 și starea FREEZE.

\par REP-14 (process hollowing) necesită activarea \texttt{kernel.yama.ptrace\_scope=0} prin sysctl, deoarece Debian 12 restricționează implicit operațiunile ptrace la relații părinte-copil. Scenariul demonstrează că Sentinel detectează corect tehnica indiferent de configurația Yama, deoarece monitorizarea se face la nivelul syscall, nu la nivelul politicii Yama.


\section{Considerente de performanță}
\label{sec:performance}

\par Penalizarea de performanță a eBPF a fost evaluată comparând utilizarea CPU pe un workload de compilare kernel (\texttt{make -j4 bzImage}) cu și fără daemonul Sentinel activ, în mașina virtuală de test. Rezultatele medii pe 5 rulări sunt prezentate în Tabelul~\ref{tab:perf}.

\begin{table}[htbp]
\begin{center}
\begin{tabular}{|p{120pt}|p{80pt}|p{80pt}|}
\hline
\textbf{Configurație} & \textbf{Durată medie (s)} & \textbf{CPU user\%} \\
\hline\hline
Fără Sentinel                & 187.4 & 94.2 \\
\hline
Cu Sentinel (fără scorer ML) & 189.7 & 95.4 \\
\hline
Cu Sentinel (cu scorer ML)   & 192.1 & 96.1 \\
\hline
\end{tabular}
\end{center}
\caption{Comparație de performanță: compilare kernel cu și fără Sentinel activ}
\label{tab:perf}
\end{table}

\par Penalizarea de performanță introdusă de daemonul eBPF fără scorer ML este de $\approx 1.2\%$ din durata totală, sub pragul nefuncțional RN1 de 2\%. Scorerul ML adaugă $\approx 2.7$ secunde suplimentare, dar acesta operează asincron și nu blochează procesarea principală; impactul reflectă consumul de CPU al inferenței Python pe thread-uri separate.

\par Latența de carantinare (de la inserarea evenimentului în baza de date la aplicarea acțiunii) a fost măsurată prin câmpul \texttt{latency\_ms} din \texttt{quarantine\_log}. Pe cele 11 scenarii din setul de bază care produc stări $\geq$ THROTTLE, latența medie înregistrată a fost de $1.8$ ms pentru THROTTLE, $3.1$ ms pentru ISOLATE (include apelurile nftables) și $0.9$ ms pentru FREEZE (\texttt{SIGSTOP} este imediat, \texttt{cgroup.freeze} este asincron). Toate valorile se încadrează în cerința RN2 ($< 5$ ms).

\par În testul REP-10 (fork bomb, 100+ apeluri \texttt{clone}/s timp de 5 secunde), nu s-au înregistrat drop-uri în ring buffer-ul de 64 MiB. Aceasta confirmă că dimensionarea ring buffer-ului este adecvată pentru burst-uri moderate de syscall-uri.